% figures/w03-thrust-to-shafts.svg — 제어기가 낸 (X, N) 이 두 축의 회전수가 되기까지.
%
%   matlab -batch "addpath _tools; w03_prop_curve"     <- 먼저 좌표를 계산하고
%   bash _tools/tikz2svg.sh figures/src/w03-thrust-to-shafts.tex figures/w03-thrust-to-shafts.svg
%
% 사용자 수업 노트(2026-10-13)의 1쪽 아래 + 2쪽을 한 그림으로. 사용자 지시:
% "어떻게 FL, FR에서 nR, nL로 가는지... n 에 대한 F 그래프도 그려주면 좋을 것 같아...
%  정회전 역회전에 대한 2차 함수 그래프"
%
% 두 단계를 **나란히** 보여 주는 것이 요점이다.
%   1 단계  선형이고 되돌릴 수 있다:  [X; N] = T [F_L; F_R],  T = [1 1; b -b]
%   2 단계  비선형이고 축마다 따로:   F = k n|n|  ->  n = +-sqrt(|F|/k)
% 학생이 "배분" 을 행렬 하나로 기억하면 2 단계를 잊는다. 그래서 한 그림에 둘 다 둔다.
%
% 좌표는 전부 _tools/w03_prop_curve.m 이 계산한다 — 이 파일에는 산술이 없다.
% 원점에서 기울기 0 인 것과 역회전이 41.8 % 약한 것이 **그림에서 그대로 보여야** 하고,
% 손으로 포물선을 찍으면 둘 다 틀린다 (CLAUDE.md 4 절 규칙 6-1).
%
% 배치 주의 (렌더해서 고친 것, 2026-10-13):
%   - 설명 상자는 **빈 사분면**에 둔다. 정회전 곡선이 1 사분면, 역회전이 3 사분면을 쓰므로
%     글은 2 사분면(왼쪽 위)과 4 사분면(오른쪽 아래)에 들어간다
%   - 축 위의 수치는 축 **바깥쪽**에. n_min 만 축 위에 둔다 — 아래는 곡선이 지나간다
%   - 회전 화살표는 선체 **밖**으로. 안에 두면 폰툰과 겹친다
\documentclass[border=6pt]{standalone}
% gnc-style.tex — 이 볼트의 TikZ 그림이 공유하는 스타일.
%
% 저널 그림 규약을 스타일 하나로 굳혀 둔다. 그림마다 선 굵기나 화살촉을
% 다시 정하지 않는다 — 그것이 그림이 제각각으로 보이는 원인이다.
%
% 쓰는 법:  % gnc-style.tex — 이 볼트의 TikZ 그림이 공유하는 스타일.
%
% 저널 그림 규약을 스타일 하나로 굳혀 둔다. 그림마다 선 굵기나 화살촉을
% 다시 정하지 않는다 — 그것이 그림이 제각각으로 보이는 원인이다.
%
% 쓰는 법:  % gnc-style.tex — 이 볼트의 TikZ 그림이 공유하는 스타일.
%
% 저널 그림 규약을 스타일 하나로 굳혀 둔다. 그림마다 선 굵기나 화살촉을
% 다시 정하지 않는다 — 그것이 그림이 제각각으로 보이는 원인이다.
%
% 쓰는 법:  \input{gnc-style}  (figures/src/ 안에서)

\usepackage{tikz}
\usepackage{amsmath}
\usepackage{xcolor}
\usetikzlibrary{arrows.meta, positioning, calc, fit, backgrounds, decorations.pathmorphing}

% ---- 색 --------------------------------------------------------------------
% 색은 강조 하나에만 쓴다. 나머지는 검정.
\definecolor{ink}{HTML}{1A1A1A}
\definecolor{accent}{HTML}{7C3AED}
\definecolor{warn}{HTML}{B91C1C}
\definecolor{soft}{HTML}{666666}

% 기하 그림(좌표계·유도)에서만 쓰는 넷. 블록선도는 위의 흑백 + accent 로 충분하다.
\definecolor{pathc}{HTML}{0D9488}   % 계획 경로 · 항적
\definecolor{errc}{HTML}{C2410C}    % 오차 (cross-track 등)
\definecolor{alng}{HTML}{B45309}    % along-track
\definecolor{flow}{HTML}{0369A1}    % 조류 · 유속

% ---- 선과 화살촉 -----------------------------------------------------------
\tikzset{
  % 신호선. 화살촉은 작고 뾰족한 것 하나뿐이다.
  sig/.style   = {ink, line width=0.5pt, -{Stealth[length=4pt, width=3pt]}},
  wire/.style  = {ink, line width=0.5pt},
  % 강조 경로 — 파선으로 구분한다. 흑백 인쇄에서도 살아야 하기 때문이다.
  asig/.style  = {accent, line width=0.5pt, densely dashed,
                  -{Stealth[length=4pt, width=3pt]}},
  awire/.style = {accent, line width=0.5pt, densely dashed},
  % 블록. 흰 바탕, 직각, 검은 테두리.
  blk/.style   = {draw=ink, line width=0.55pt, fill=white,
                  minimum width=17mm, minimum height=8mm, inner sep=2pt, font=\small},
  % 합산점
  sum/.style   = {draw=ink, line width=0.55pt, fill=white, circle,
                  inner sep=0pt, minimum size=4.6mm, font=\scriptsize},
  asum/.style  = {draw=accent, line width=0.55pt, fill=white, circle,
                  inner sep=0pt, minimum size=4.6mm, font=\scriptsize},
  % 분기점
  dot/.style   = {ink, fill=ink, circle, inner sep=0pt, minimum size=1.5mm},
  % 신호 이름 — 선 위에 작게
  sname/.style = {font=\small, inner sep=1.5pt},
  % 그림 안의 짧은 설명. 그림 안의 글은 최소로 둔다
  note/.style  = {font=\scriptsize, soft, inner sep=1.5pt},
  anote/.style = {font=\scriptsize, accent, inner sep=1.5pt},
  % 축
  axis/.style  = {ink, line width=0.5pt},
  tick/.style  = {font=\scriptsize, soft},
}

% 캡션 한 줄 — 그림 아래 가는 선과 함께
\newcommand{\gnccaption}[2]{%
  \node[anchor=north west, text width=#1, align=left, font=\scriptsize, ink]
        at (current bounding box.south west) {\vspace{2pt}#2};}
  (figures/src/ 안에서)

\usepackage{tikz}
\usepackage{amsmath}
\usepackage{xcolor}
\usetikzlibrary{arrows.meta, positioning, calc, fit, backgrounds, decorations.pathmorphing}

% ---- 색 --------------------------------------------------------------------
% 색은 강조 하나에만 쓴다. 나머지는 검정.
\definecolor{ink}{HTML}{1A1A1A}
\definecolor{accent}{HTML}{7C3AED}
\definecolor{warn}{HTML}{B91C1C}
\definecolor{soft}{HTML}{666666}

% 기하 그림(좌표계·유도)에서만 쓰는 넷. 블록선도는 위의 흑백 + accent 로 충분하다.
\definecolor{pathc}{HTML}{0D9488}   % 계획 경로 · 항적
\definecolor{errc}{HTML}{C2410C}    % 오차 (cross-track 등)
\definecolor{alng}{HTML}{B45309}    % along-track
\definecolor{flow}{HTML}{0369A1}    % 조류 · 유속

% ---- 선과 화살촉 -----------------------------------------------------------
\tikzset{
  % 신호선. 화살촉은 작고 뾰족한 것 하나뿐이다.
  sig/.style   = {ink, line width=0.5pt, -{Stealth[length=4pt, width=3pt]}},
  wire/.style  = {ink, line width=0.5pt},
  % 강조 경로 — 파선으로 구분한다. 흑백 인쇄에서도 살아야 하기 때문이다.
  asig/.style  = {accent, line width=0.5pt, densely dashed,
                  -{Stealth[length=4pt, width=3pt]}},
  awire/.style = {accent, line width=0.5pt, densely dashed},
  % 블록. 흰 바탕, 직각, 검은 테두리.
  blk/.style   = {draw=ink, line width=0.55pt, fill=white,
                  minimum width=17mm, minimum height=8mm, inner sep=2pt, font=\small},
  % 합산점
  sum/.style   = {draw=ink, line width=0.55pt, fill=white, circle,
                  inner sep=0pt, minimum size=4.6mm, font=\scriptsize},
  asum/.style  = {draw=accent, line width=0.55pt, fill=white, circle,
                  inner sep=0pt, minimum size=4.6mm, font=\scriptsize},
  % 분기점
  dot/.style   = {ink, fill=ink, circle, inner sep=0pt, minimum size=1.5mm},
  % 신호 이름 — 선 위에 작게
  sname/.style = {font=\small, inner sep=1.5pt},
  % 그림 안의 짧은 설명. 그림 안의 글은 최소로 둔다
  note/.style  = {font=\scriptsize, soft, inner sep=1.5pt},
  anote/.style = {font=\scriptsize, accent, inner sep=1.5pt},
  % 축
  axis/.style  = {ink, line width=0.5pt},
  tick/.style  = {font=\scriptsize, soft},
}

% 캡션 한 줄 — 그림 아래 가는 선과 함께
\newcommand{\gnccaption}[2]{%
  \node[anchor=north west, text width=#1, align=left, font=\scriptsize, ink]
        at (current bounding box.south west) {\vspace{2pt}#2};}
  (figures/src/ 안에서)

\usepackage{tikz}
\usepackage{amsmath}
\usepackage{xcolor}
\usetikzlibrary{arrows.meta, positioning, calc, fit, backgrounds, decorations.pathmorphing}

% ---- 색 --------------------------------------------------------------------
% 색은 강조 하나에만 쓴다. 나머지는 검정.
\definecolor{ink}{HTML}{1A1A1A}
\definecolor{accent}{HTML}{7C3AED}
\definecolor{warn}{HTML}{B91C1C}
\definecolor{soft}{HTML}{666666}

% 기하 그림(좌표계·유도)에서만 쓰는 넷. 블록선도는 위의 흑백 + accent 로 충분하다.
\definecolor{pathc}{HTML}{0D9488}   % 계획 경로 · 항적
\definecolor{errc}{HTML}{C2410C}    % 오차 (cross-track 등)
\definecolor{alng}{HTML}{B45309}    % along-track
\definecolor{flow}{HTML}{0369A1}    % 조류 · 유속

% ---- 선과 화살촉 -----------------------------------------------------------
\tikzset{
  % 신호선. 화살촉은 작고 뾰족한 것 하나뿐이다.
  sig/.style   = {ink, line width=0.5pt, -{Stealth[length=4pt, width=3pt]}},
  wire/.style  = {ink, line width=0.5pt},
  % 강조 경로 — 파선으로 구분한다. 흑백 인쇄에서도 살아야 하기 때문이다.
  asig/.style  = {accent, line width=0.5pt, densely dashed,
                  -{Stealth[length=4pt, width=3pt]}},
  awire/.style = {accent, line width=0.5pt, densely dashed},
  % 블록. 흰 바탕, 직각, 검은 테두리.
  blk/.style   = {draw=ink, line width=0.55pt, fill=white,
                  minimum width=17mm, minimum height=8mm, inner sep=2pt, font=\small},
  % 합산점
  sum/.style   = {draw=ink, line width=0.55pt, fill=white, circle,
                  inner sep=0pt, minimum size=4.6mm, font=\scriptsize},
  asum/.style  = {draw=accent, line width=0.55pt, fill=white, circle,
                  inner sep=0pt, minimum size=4.6mm, font=\scriptsize},
  % 분기점
  dot/.style   = {ink, fill=ink, circle, inner sep=0pt, minimum size=1.5mm},
  % 신호 이름 — 선 위에 작게
  sname/.style = {font=\small, inner sep=1.5pt},
  % 그림 안의 짧은 설명. 그림 안의 글은 최소로 둔다
  note/.style  = {font=\scriptsize, soft, inner sep=1.5pt},
  anote/.style = {font=\scriptsize, accent, inner sep=1.5pt},
  % 축
  axis/.style  = {ink, line width=0.5pt},
  tick/.style  = {font=\scriptsize, soft},
}

% 캡션 한 줄 — 그림 아래 가는 선과 함께
\newcommand{\gnccaption}[2]{%
  \node[anchor=north west, text width=#1, align=left, font=\scriptsize, ink]
        at (current bounding box.south west) {\vspace{2pt}#2};}

% 자동 생성 — w03_prop_curve.m 가 만든다. 손으로 고치지 않는다.
% generated by _tools/w03_prop_curve.m; do not edit by hand
\def\kpos{0.01108}
\def\kneg{0.00645}
\def\nmax{103.9}
\def\nmin{-101.7}
\def\Fmax{119.7}
\def\Fmin{-66.7}
\def\ypont{0.3950}
\def\axL{-6.0526}
\def\axR{6.2105}
\def\axD{-3.0769}
\def\axU{5.1923}
\def\nmaxx{5.4700}
\def\nminx{-5.3546}
\def\Fmaxy{4.6032}
\def\Fminy{-2.5657}
\def\Fw{38.16}
\def\nw{58.69}
\def\Fwy{1.4677}
\def\nwx{3.0887}
\def\curveAhead{(0.0000,0.0000) -- (0.0927,0.0013) -- (0.1854,0.0053) -- (0.2781,0.0119) -- (0.3709,0.0212) -- (0.4636,0.0331) -- (0.5563,0.0476) -- (0.6490,0.0648) -- (0.7417,0.0846) -- (0.8344,0.1071) -- (0.9271,0.1322) -- (1.0198,0.1600) -- (1.1126,0.1904) -- (1.2053,0.2235) -- (1.2980,0.2592) -- (1.3907,0.2975) -- (1.4834,0.3385) -- (1.5761,0.3822) -- (1.6688,0.4284) -- (1.7615,0.4774) -- (1.8543,0.5289) -- (1.9470,0.5832) -- (2.0397,0.6400) -- (2.1324,0.6995) -- (2.2251,0.7617) -- (2.3178,0.8265) -- (2.4105,0.8939) -- (2.5032,0.9640) -- (2.5960,1.0367) -- (2.6887,1.1121) -- (2.7814,1.1901) -- (2.8741,1.2708) -- (2.9668,1.3541) -- (3.0595,1.4401) -- (3.1522,1.5287) -- (3.2449,1.6199) -- (3.3377,1.7138) -- (3.4304,1.8103) -- (3.5231,1.9095) -- (3.6158,2.0113) -- (3.7085,2.1158) -- (3.8012,2.2229) -- (3.8939,2.3327) -- (3.9866,2.4451) -- (4.0794,2.5601) -- (4.1721,2.6778) -- (4.2648,2.7981) -- (4.3575,2.9211) -- (4.4502,3.0467) -- (4.5429,3.1750) -- (4.6356,3.3059) -- (4.7283,3.4395) -- (4.8211,3.5757) -- (4.9138,3.7145) -- (5.0065,3.8560) -- (5.0992,4.0002) -- (5.1919,4.1469) -- (5.2846,4.2964) -- (5.3773,4.4484) -- (5.4700,4.6032)}
\def\curveAstern{(-5.3546,-2.5657) -- (-5.2638,-2.4795) -- (-5.1731,-2.3947) -- (-5.0823,-2.3114) -- (-4.9915,-2.2296) -- (-4.9008,-2.1493) -- (-4.8100,-2.0704) -- (-4.7193,-1.9930) -- (-4.6285,-1.9171) -- (-4.5378,-1.8426) -- (-4.4470,-1.7697) -- (-4.3563,-1.6982) -- (-4.2655,-1.6282) -- (-4.1747,-1.5596) -- (-4.0840,-1.4925) -- (-3.9932,-1.4269) -- (-3.9025,-1.3628) -- (-3.8117,-1.3002) -- (-3.7210,-1.2390) -- (-3.6302,-1.1793) -- (-3.5395,-1.1211) -- (-3.4487,-1.0643) -- (-3.3579,-1.0090) -- (-3.2672,-0.9552) -- (-3.1764,-0.9029) -- (-3.0857,-0.8520) -- (-2.9949,-0.8027) -- (-2.9042,-0.7547) -- (-2.8134,-0.7083) -- (-2.7227,-0.6634) -- (-2.6319,-0.6199) -- (-2.5411,-0.5779) -- (-2.4504,-0.5373) -- (-2.3596,-0.4982) -- (-2.2689,-0.4607) -- (-2.1781,-0.4245) -- (-2.0874,-0.3899) -- (-1.9966,-0.3567) -- (-1.9059,-0.3250) -- (-1.8151,-0.2948) -- (-1.7244,-0.2661) -- (-1.6336,-0.2388) -- (-1.5428,-0.2130) -- (-1.4521,-0.1887) -- (-1.3613,-0.1658) -- (-1.2706,-0.1445) -- (-1.1798,-0.1246) -- (-1.0891,-0.1061) -- (-0.9983,-0.0892) -- (-0.9076,-0.0737) -- (-0.8168,-0.0597) -- (-0.7260,-0.0472) -- (-0.6353,-0.0361) -- (-0.5445,-0.0265) -- (-0.4538,-0.0184) -- (-0.3630,-0.0118) -- (-0.2723,-0.0066) -- (-0.1815,-0.0029) -- (-0.0908,-0.0007) -- (0.0000,0.0000)}


\definecolor{knownc}{HTML}{1D4ED8}

\tikzset{
  hull/.style    = {ink, line width=0.9pt, fill=black!4},
  pont/.style    = {ink, line width=0.7pt, fill=pathc!12},
  thrust/.style  = {errc, line width=1.4pt, -{Stealth[length=5.0pt,width=3.6pt]}},
  armline/.style = {knownc, line width=0.8pt,
                    {Stealth[length=3.4pt,width=2.4pt]}-{Stealth[length=3.4pt,width=2.4pt]}},
  ctr/.style     = {soft, line width=0.4pt, densely dashed},
  spin/.style    = {alng, line width=1.1pt, -{Stealth[length=4.4pt,width=3.2pt]}},
  ax/.style      = {ink, line width=0.7pt, -{Stealth[length=4.6pt,width=3.2pt]}},
  tick/.style    = {soft, line width=0.3pt, densely dotted},
  curveA/.style  = {errc, line width=1.6pt},
  curveB/.style  = {flow, line width=1.6pt},
  guide/.style   = {accent, line width=0.8pt, densely dashed},
  lead/.style    = {soft, line width=0.35pt, -{Stealth[length=3.2pt,width=2.2pt]}},
  lbl/.style     = {font=\scriptsize},
  known/.style   = {draw=knownc, line width=0.6pt, fill=white, inner sep=3.2pt,
                    rounded corners=1pt, font=\small},
  want/.style    = {draw=accent, line width=0.9pt, fill=white, inner sep=3.4pt,
                    rounded corners=1pt, font=\small},
  capt/.style    = {font=\scriptsize\bfseries, warn},
  stage/.style   = {font=\small\bfseries, ink},
}

\begin{document}
\begin{tikzpicture}[x=1cm, y=1cm]

% =====================================================================
% 윗줄 왼쪽 — 자유물체도. 두 추력이 어디에 붙어 있는가
% =====================================================================
\begin{scope}[shift={(0,0)}]
  \node[stage, anchor=west] at (-1.9,3.55) {STEP 1 \ \ two thrusts into two forces, and back};

  \draw[ctr] (0,-2.0) -- (0,2.6);
  \draw[pont] (-1.55,-1.5) rectangle (-0.85,1.1);
  \draw[pont] ( 0.85,-1.5) rectangle ( 1.55,1.1);
  \draw[hull] (-1.55,0.25) rectangle (1.55,1.1);

  \draw[ax] (0,0) -- (0,2.30);   \node[lbl, anchor=west]  at (0.08,2.22) {$x_b$};
  \draw[ax] (0,0) -- (2.30,0);   \node[lbl, anchor=south] at (2.20,0.07) {$y_b$};
  \fill[ink] (0,0) circle (1.2pt);

  \draw[armline] (-1.2,-1.0) -- (0,-1.0);
  \draw[armline] ( 0,-1.0)   -- (1.2,-1.0);
  \node[lbl, knownc, anchor=north] at (-0.6,-1.08) {$b$};
  \node[lbl, knownc, anchor=north] at ( 0.6,-1.08) {$b$};
  \node[lbl, knownc, anchor=north] at (0,-1.62) {$b = \ypont$ m};

  \draw[thrust] (-1.2,1.1) -- (-1.2,2.60);
  \draw[thrust] ( 1.2,1.1) -- ( 1.2,2.00);
  \node[lbl, errc, anchor=south] at (-1.2,2.63) {$F_L$};
  \node[lbl, errc, anchor=south] at ( 1.2,2.03) {$F_R$};

  %  회전 화살표는 선체 밖으로 / the turn arrow goes outside the hull
  \draw[spin] (1.00,-1.73) arc (-60:-18:2.0);
  \node[lbl, alng, anchor=west] at (2.00,-0.78) {$N>0$};

  \node[capt, anchor=north] at (0.2,-2.25) {where the two thrusts act};
\end{scope}

% =====================================================================
% 윗줄 오른쪽 — 행렬과 그 역
% =====================================================================
\begin{scope}[shift={(4.9,0)}]
  \node[known, anchor=north west, align=left] at (0,3.10)
    {both push forward, and the two arms are opposite\\[3pt]
     $\begin{bmatrix} X \\ N \end{bmatrix}
      = \underbrace{\begin{bmatrix} 1 & 1 \\ b & -b \end{bmatrix}}_{\mathbf{T}}
        \begin{bmatrix} F_L \\ F_R \end{bmatrix}$};
  \node[anchor=north west, align=left, font=\scriptsize] at (0,1.20)
    {$\det\mathbf{T} = -2b \neq 0$: the matrix is invertible, so the split is \emph{unique}};
  \node[want, anchor=north west, align=left] at (0,0.70)
    {$\begin{bmatrix} F_L \\ F_R \end{bmatrix}
      = \mathbf{T}^{-1}\begin{bmatrix} X \\ N \end{bmatrix}
      = \begin{bmatrix} \tfrac12 & \tfrac{1}{2b} \\[2pt] \tfrac12 & -\tfrac{1}{2b} \end{bmatrix}
        \begin{bmatrix} X \\ N \end{bmatrix}$};
  \node[anchor=north west, align=left, font=\scriptsize] at (0,-1.25)
    {each propeller carries \emph{half} the surge force, and the moment\\
     is added to one and taken from the other};
  \node[capt, anchor=north] at (2.7,-2.25) {linear, and exactly invertible};
\end{scope}

% =====================================================================
% 아랫줄 — 추력 곡선 F = k n|n| 과 그 역산
% =====================================================================
\begin{scope}[shift={(2.3,-9.4)}]
  \node[stage, anchor=west] at (\axL,5.95)
    {STEP 2 \ \ each force onto its own shaft --- nonlinear, one propeller at a time};

  % 축 / the axes
  %  축 이름은 축 **위**, 눈금 수치는 축 아래 — 같은 쪽에 두면 겹친다 (렌더해서 고침)
  \draw[ax] (\axL,0) -- (\axR,0);  \node[lbl, anchor=south east] at (\axR,0.14) {$n$ \ [rad/s]};
  \draw[ax] (0,\axD) -- (0,\axU);  \node[lbl, anchor=south west] at (0.10,\axU) {$F$ \ [N]};

  % 곡선 두 가지 / the two branches
  \draw[curveB] \curveAstern;
  \draw[curveA] \curveAhead;

  % 한계 / the limits, labelled outside the axes
  \draw[tick] (\nmaxx,0) -- (\nmaxx,\Fmaxy) -- (0,\Fmaxy);
  \draw[tick] (\nminx,0) -- (\nminx,\Fminy) -- (0,\Fminy);
  \fill[errc] (\nmaxx,\Fmaxy) circle (1.7pt);
  \fill[flow] (\nminx,\Fminy) circle (1.7pt);
  \node[lbl, errc, anchor=east]  at (-0.12,\Fmaxy) {$F_{\max}=\Fmax$};
  \node[lbl, flow, anchor=east]  at (-0.12,\Fminy) {$F_{\min}=\Fmin$};
  \node[lbl, errc, anchor=north] at (\nmaxx,-0.14) {$n_{\max}=\nmax$};
  \node[lbl, flow, anchor=south] at (\nminx, 0.14) {$n_{\min}=\nmin$};

  % 2 사분면 (왼쪽 위) 은 비어 있다 — 설명을 여기 둔다
  \node[known, anchor=north west, align=left] at (\axL,\axU)
    {\textbf{ahead}\ \ $F = k_{pos}\, n\lvert n\rvert$\\[2pt] $k_{pos} = \kpos$};
  \node[anchor=north west, align=left, font=\scriptsize, soft] at (\axL,2.55)
    {\textbf{flat at the origin.}\\
     $\mathrm{d}F/\mathrm{d}n = 0$ there, so a\\
     slow shaft makes almost no\\
     thrust, and a small $F$ needs\\
     a surprisingly large step in $n$};
  \draw[lead] (-3.35,1.00) .. controls (-1.6,0.55) .. (-0.18,0.06);

  % 3 사분면 아래 — 역회전 쪽 설명 / the astern branch, below its own curve
  \node[anchor=north west, align=left, font=\scriptsize, flow] at (\axL,-2.70)
    {\textbf{astern}\ \ $F = k_{neg}\, n\lvert n\rvert$, \ $k_{neg} = \kneg$\\[2pt]
     the same blade turned backwards: \textbf{41.8\,\% weaker},\\
     which is why $\lvert F_{\min}\rvert < F_{\max}$};

  % 4 사분면 (오른쪽 아래) 은 비어 있다 — 역산을 여기 둔다
  \node[want, anchor=north west, align=left] at (1.25,-0.75)
    {$n = \begin{cases}
        +\sqrt{F/k_{pos}}, & F \ge 0\\[3pt]
        -\sqrt{\lvert F\rvert/k_{neg}}, & F < 0
      \end{cases}$};
  \node[anchor=north west, align=left, font=\scriptsize] at (1.25,-2.45)
    {read the curve \emph{backwards}. Each propeller is read on its own,\\
     because each has its own $F$ and neither can help the other};

  % 읽어 내는 예 / the worked reading
  \draw[guide] (0,\Fwy) -- (\nwx,\Fwy) -- (\nwx,0);
  \fill[accent] (\nwx,\Fwy) circle (1.7pt);
  \node[lbl, accent, anchor=south west] at (0.15,\Fwy+0.06) {$F_L = \Fw$};
  \node[lbl, accent, anchor=north] at (\nwx,-0.14) {$n_L = \nw$};

  \node[capt, anchor=north] at (0,-4.10)
    {nonlinear, but still one-to-one: every reachable $F$ has exactly one $n$};
\end{scope}

\end{tikzpicture}
\end{document}
